\subsection{Full-matrix presentation of a factor with scalar centre}

\begin{theorem}[Scalar centre forces a full matrix algebra of a definite size]
    \label{thm:qca_central_star_subalgebra_matrix}
    \lean{StarSubalgebra.exists_starAlgEquiv_matrix_and_finrank_of_isCentral}
    \leanok
    Let $I$ be a finite nonempty set and let $S\subseteq M_I(\C)$ be a unital
    star-subalgebra whose centre consists of the multiples of its unit,
    \begin{align}
        \{z\in S\mid zs=sz\text{ for every }s\in S\}
         & =\C\mathbf 1.
        \label{eq:qca_central_star_subalgebra_scalar_centre}
    \end{align}
    Then there is an $r\geq1$ with
    \begin{align}
        S            & \cong\MN{r}, \notag \\
        \dim_{\C}(S) & =r^2,
    \end{align}
    the first isomorphism being one of star-algebras.

    Nonemptiness of $I$ cannot be dropped.  For $I=\emptyset$ the ambient
    algebra $M_I(\C)$ is the zero algebra, its unique subalgebra still
    satisfies (\ref{eq:qca_central_star_subalgebra_scalar_centre}), and it is
    isomorphic to no $\MN{r}$ with $r\geq1$.

    \textbf{Scope restriction (Schumacher--Werner \texttt{Csform}):}
    This is the single building block of the structure theorem for
    finite-dimensional $C^*$-algebras, Proposition~\texttt{Csform}
    \cite[lines~2082--2098]{SchumacherWerner2004QCA}, for an algebra already
    presented inside a matrix algebra.  It asserts neither the direct-sum
    decomposition $\mathcal A\cong\bigoplus_\mu\MN{n_\mu}$ of that
    proposition, nor its cut by the minimal projections of the centre, nor the
    abstract $C^*$-algebra generality in which it is stated.  The omitted
    content and its elimination plan are recorded in
    \cite{gap:sw04_csform_matrix_scope}.
\end{theorem}

\begin{proof}\leanok
    A star-subalgebra of a complex matrix algebra is semisimple, so it is
    isomorphic to a finite product $\prod_{k=1}^{m}\MN{d_k}$ with every
    $d_k\geq1$.  The element that is the unit in the $k$-th coordinate and
    zero in the others is central, hence by
    (\ref{eq:qca_central_star_subalgebra_scalar_centre}) a multiple
    $\lambda\mathbf 1$; comparing its $k$-th coordinate with another one gives
    $\lambda=1$ and $\lambda=0$, so $m\leq1$.  The unit of $M_I(\C)$ is
    nonzero because $I$ is nonempty, so $m=1$ and $S$ is a simple ring.  The
    Wedderburn theorem over the algebraically closed field $\C$ then produces
    an $r\geq1$ and an isomorphism of $\C$-algebras $S\cong\MN{r}$, and
    conjugating it by a positive definite matrix makes it preserve adjoints
    \cite[``A $*$-subalgebra with scalar centre is a full matrix
        algebra'']{QICLean2026Blueprint}.  A star-isomorphism is in particular
    a $\C$-linear isomorphism, and $\dim_{\C}\MN{r}=r^2$
    \cite[``Dimension of a subalgebra presented as a full matrix
        algebra'']{QICLean2026Blueprint}.
\end{proof}

Gross, Nesme, Vogts and Werner apply this conclusion to the site-indexed
support algebras \cite[lines~1276--1282]{Gross2012IndexTheory}: an element of
$\mathcal R_x$ commuting with all of $\mathcal R_x$ commutes with every
$\mathcal R_y$, hence with the algebra they generate, which is the whole
quasi-local algebra; its centre is scalar, so
(\ref{eq:qca_central_star_subalgebra_scalar_centre}) holds for $\mathcal R_x$
and $\mathcal R_x\cong M_{r(x)}(\C)$.
